\documentclass[11pt,a4paper]{scrartcl}

\usepackage{../manostilius_lt}

\title{Vienanaris}
\author{Andrius Berniukevičius}

\begin{document}

\maketitle

\section{Vienanario apibrėžimas}

Algebriniai reiškiniai gali būti sudaryti iš skaičių, raidžių ir įvairių veiksmų. Vienas paprasčiausių algebrinių reiškinių yra vienanaris.

\begin{definition}
\vocab{Vienanaris} -- tai skaičiaus ir kintamųjų, pakeltų natūraliaisiais laipsniais, sandauga.

Vienanaryje neatliekami sudėties arba atimties veiksmai.
\end{definition}

Vienanarių pavyzdžiai:
\[
5a, \qquad -3x^2y, \qquad \frac{2}{7}ab^3, \qquad 8.
\]

\begin{example}
Nustatykime, kurie reiškiniai yra vienanariai:

\begin{tasks}[style=itemize](2)
    \task $7x^3y$
    \task $4a-3b$
    \task $-\frac{5}{2}m^2n^4$
    \task $x^2+5$
    \task $\frac{3a}{5}$
    \task $(a+b)^2$
    \task $4x^{-2}y$
    \task $3a\sqrt{b}$
\end{tasks}
\textbf{Sprendimas}\\
Vienanariai yra:
\[
7x^3y, \qquad -\frac{5}{2}m^2n^4, \qquad \frac{3a}{5}.
\]
Reiškiniai $4a-3b$, $x^2+5$ ir $(a+b)^2$ nėra vienanariai, nes juose atliekamas sudėties arba atimties veiksmas. Reiškinys $4x^{-2}y$ nėra vienanaris, nes kintamojo $x$ rodiklis yra neigiamas. Reiškinys $3a\sqrt{b}$ taip pat nėra vienanaris, nes kintamasis $b$ yra po šaknies ženklu.
\end{example}

\begin{remark}
Skaičius taip pat laikomas vienanariu. Pavyzdžiui:
\[
8=8 \cdot 1.
\]

Trupmena laikoma vienanariu, jei vardiklyje yra tik skaičius:
\[
\frac{3a}{5}=\frac{3}{5}\cdot a.
\]
\end{remark}

\begin{remark}
Reiškinys su kintamuoju vardiklyje nėra vienanaris. Pavyzdžiui:
\[
\frac{5}{x}
\]
nėra vienanaris, nes jį galima užrašyti kaip $5x^{-1}$, o rodiklis $-1$ nėra natūralusis skaičius.
\end{remark}

\section{Vienanario skaitinė ir raidinė dalys}

Kiekvieną vienanarį galime išskaidyti į skaitinę ir raidinę dalis.

\begin{definition}
\vocab{Vienanario koeficientas} -- tai skaitinis vienanario daugiklis.
\end{definition}

\begin{definition}
\vocab{Vienanario raidinė dalis} -- tai vienanaryje esančių kintamųjų ir jų laipsnių sandauga.
\end{definition}

\begin{example}
Panagrinėkime vienanarį:
\[
-6a^3b^2.
\]
\textbf{Sprendimas}\\
Jo koeficientas yra $-6$, o raidinė dalis yra:
\[
a^3b^2.
\]
\end{example}

\begin{example}
Panagrinėkime vienanarį:
\[
x^4y.
\]
\textbf{Sprendimas}\\
Jo koeficientas yra $1$, nes:
\[
x^4y=1 \cdot x^4y.
\]
Raidinė dalis yra $x^4y$.
\end{example}

\begin{example}
Panagrinėkime vienanarį:
\[
-abc^2.
\]
\textbf{Sprendimas}\\
Jo koeficientas yra $-1$, nes:
\[
-abc^2=-1 \cdot abc^2.
\]
Raidinė dalis yra $abc^2$.
\end{example}

\section{Vienanario laipsnis}

Vienanario laipsnis priklauso tik nuo raidinės dalies.

\begin{definition}
\vocab{Vienanario laipsnis} -- tai visų vienanario raidinėje dalyje esančių kintamųjų laipsnių rodiklių suma.
Nelygaus nuliui skaičiaus laipsnis yra $0$.
\end{definition}

\begin{example}
Raskime vienanario $5a^2b^3$ laipsnį.\\ \\
\textbf{Sprendimas}\\
Kintamojo $a$ rodiklis yra $2$, o kintamojo $b$ rodiklis yra $3$. Todėl:
\[
2+3=5.
\]
Vienanario $5a^2b^3$ laipsnis yra $5$.
\end{example}

\begin{example}
Raskime vienanario $-7x^4y^2z$ laipsnį.\\ \\
\textbf{Sprendimas}\\
Kintamojo $z$ rodiklis yra $1$, todėl:
\[
4+2+1=7.
\]
Vienanario $-7x^4y^2z$ laipsnis yra $7$.
\end{example}

\begin{example}
Raskime vienanario $12$ laipsnį. \\ \\
\textbf{Sprendimas}\\ 
Kadangi skaičius neturi raidinės dalies, jo laipsnis yra:
\[
0.
\]
\end{example}

\begin{remark}
Jeigu prie raidės laipsnio rodiklis neparašytas, jis yra lygus $1$:
\[
a=a^1, \qquad xy=x^1y^1.
\]
Pavyzdžiui, vienanario $4x^2y$ laipsnis yra:
\[
2+1=3.
\]
\end{remark}
\newpage
\section{Uždaviniai}

\begin{problem}
Nustatykite, kurie reiškiniai yra vienanariai.

\begin{tasks}[label={(\arabic*)}, label-offset=0.9em](2)
    \task $8a^2b$
    \task $3x-5$
    \task $-\frac{4}{7}m^3n$
    \task $a^2+b^2$
    \task $36$
    \task $\frac{x+3}{5}$
    \task $-xyz^4$
    \task $\frac{5}{a}$
    \task $0,6p^2q^3$
    \task $(x-2)(x+2)$
    \task $\frac{7a^2}{9}$
    \task $-2a^{-3}b^4$
    \task $-\frac{5x^3y}{12}$
    \task $5x\sqrt{y}$
    \task $\frac{3}{8}p^2q^4$
    \task $\frac{4m^3}{n}$
    \task $\frac{3}{4}m^2\sqrt[3]{n}$
    \task $7pq^{-1}$
\end{tasks}
\end{problem}

\begin{problem}
Nurodykite kiekvieno vienanario koeficientą ir raidinę dalį.

\begin{tasks}[label={(\arabic*)}, label-offset=0.9em](2)
    \task $7a^3b$
    \task $-5x^2y^4$
    \task $mn^2$
    \task $\frac{7a^2b^3}{9}$
    \task $-\frac{3}{8}p^5q$
    \task $12$
    \task $-abc^3$
    \task $0,4x^2y$
    \task $\frac{5}{6}a^4b^2$
    \task $-\frac{11m^4n}{5}$
\end{tasks}
\end{problem}

\begin{problem}
Raskite vienanarių laipsnius.

\begin{tasks}[label={(\arabic*)}, label-offset=0.9em](2)
    \task $4a^3$
    \task $-2x^2y^5$
    \task $7abc$
    \task $-5m^4n^2p^3$
    \task $13$
    \task $\frac{1}{2}x^6y$
    \task $-a^2b^2c^2$
    \task $0,8k^3l^4$
\end{tasks}
\end{problem}

\newpage
\answerssection

\begin{problem} \phantom{text}
\begin{tasks}[label={(\arabic*)}, label-offset=0.9em](2)
    \task vienanaris;
    \task ne vienanaris;
    \task vienanaris;
    \task ne vienanaris;
    \task vienanaris;
    \task ne vienanaris;
    \task vienanaris;
    \task ne vienanaris;
    \task vienanaris;
    \task ne vienanaris.
    \task vienanaris;
    \task ne vienanaris;
    \task vienanaris;
    \task ne vienanaris;
    \task vienanaris;
    \task ne vienanaris;
    \task ne vienanaris;
    \task ne vienanaris.
\end{tasks}
\end{problem}

\begin{problem} \phantom{text}
    \begin{tasks}[label={(\arabic*)}, label-offset=0.9em](2)
    \task koeficientas $7$, raidinė dalis $a^3b$;
    \task koeficientas $-5$, raidinė dalis $x^2y^4$;
    \task koeficientas $1$, raidinė dalis $mn^2$;
    \task koeficientas $\frac{7}{9}$, raidinė dalis $a^2b^3$;
    \task koeficientas $-\frac{3}{8}$, raidinė dalis $p^5q$;
    \task koeficientas $12$, raidinės dalies nėra;
    \task koeficientas $-1$, raidinė dalis $abc^3$;
    \task koeficientas $0,4$, raidinė dalis $x^2y$;
    \task koeficientas $\frac{5}{6}$, raidinė dalis $a^4b^2$.
    \task koeficientas $-\frac{11}{5}$, raidinė dalis $m^4n$.
\end{tasks}
\end{problem}

\begin{problem} \phantom{text}
\begin{tasks}[label={(\arabic*)}, label-offset=0.9em](2)
    \task $3$;
    \task $7$;
    \task $3$;
    \task $9$;
    \task $0$;
    \task $7$;
    \task $6$;
    \task $7$.
\end{tasks}
\end{problem}

\end{document}